% ---------------------------------------------------------------------
% 第 6 章 仿真验证
% ---------------------------------------------------------------------
\section{仿真验证}

\subsection{平台与任务}

在 CoppeliaSim 中以 7-DoF KUKA LBR4+ 力矩模式仿真（控制周期 5 ms），名义动力学取自 \cite{Gaz14}。$t=2.5\,\text{s}$ 时 0.25 kg 未建模负载刚性附着，构成持续偏差型扰动 $d_{\mathrm{ex}}$。任务为抓取--搬运--圆周跟踪（半径 0.06 m，标准角速度 1.0 rad/s，高速工况 2.5 rad/s）。对比三种控制器：\textbf{C1}（本文，式 \eqref{eq:5.2}）、\textbf{C2}（\cite{Ch20}，twist 误差与本文相同但位姿反馈取螺旋对数）、\textbf{C3}（\cite{P2} 经速度伺服桥接至加速度级）。三律共享同一轨迹、初始条件与名义模型，仅 $\ddot{\boldsymbol q}_{\mathrm{ref}}$ 计算分支不同；DC 刚度均配平至 80，闭环极点 $\{-4,-20\}$，以隔离结构差异。

\subsection{验证目标：从定理到可观测预言}

仿真不追求“性能竞赛”式的百分比优胜，而是检验 §3--5 理论框架的三类可证伪预言：

\textbf{(V1) 静态刚度标度律 \eqref{eq:5.9}}：若定理 3(d) 的扰动模型与 §5.4 的线性化正确，则同一物理扰动 $d_{\mathrm{ex}}$ 应由不同增益档的稳态残差反演出一致幅值；且 1/4 旋转刚度折减因子应被独立确认。

\textbf{(V2) 证书分化的极限暴露}：C1 与 C2 共享解析前馈与几何一致 twist 误差，预期在近恒等域数值等价（确认“同信息集”）；但 C1 具备定理 3 的严格耗散等式与 H∞ 证书，C2/C3 没有。这一分化应在离开理想条件的极限工况（高速、噪声）下表现为 C1/C2 与 C3 的结构性分离——而非随机波动。

\textbf{(V3) 水平集工程余度}：定理 3(b) 给出水平集 $\Omega_c$ 的正向不变性，预期在实际运行中应留有充分余度，以确认局部结论的工程可用性。

\subsection{结果}

\subsubsection{静态刚度标度律与扰动反演（验证 V1）}

带载稳态（$t\ge 12.5\,\text{s}$）下，C1-tuned 与 C1-fast 的位姿残差如下：

\begin{center}
\begin{tabular}{@{}lccc@{}}
\toprule
档位 & $K_p=\mathrm{diag}(p_O I_3, p_T I_3)$ & $\|\mathcal T\|_{\mathrm{rms}}$ (m) & $\|\mathcal O\|_{\mathrm{rms}}$ \\
\midrule
tuned & $(320, 80)$ & $4.86\times10^{-3}$ & $4.27\times10^{-3}$ \\
fast  & $(720, 180)$ & $2.20\times10^{-3}$ & $1.93\times10^{-3}$ \\
\bottomrule
\end{tabular}
\end{center}

由 \eqref{eq:5.9} 反演等效扰动：
\begin{equation*}
\|d_v\| = k_{p,T}\|\mathcal T\|_{\mathrm{ss}}, \qquad \|d_\omega\| = \tfrac12\lambda(K_{p,O})\|\mathcal O\|_{\mathrm{ss}}.
\end{equation*}
得 tuned 档 $(\|d_v\|, \|d_\omega\|) = (0.389, 0.683)$，fast 档 $(0.396, 0.696)$，两档一致至 \textbf{1.93\%}。该强一致性（两个独立增益档反演出同一物理量）直接验证了定理 3(d) 扰动模型的结构正确性——若 1/4 旋转折减因子缺失，旋转反演值将与平移分量相差 3.5 倍，一致性立即破坏。

作为对照，base 档（$K_p=16I_6$，未补偿 1/4 折减，有效旋转刚度仅 4）带载姿态误差放大至 tuned 档的 12.3 倍，反演扰动偏离 35\%--38\%。这从反面验证了 §5.4 的 1/4 折减补偿规则的必要性：不补偿则旋转/平移带宽严重失配，\eqref{eq:5.9} 失效。

\subsubsection{极限工况下的证书分化（验证 V2）}

标准工况下，三律位置级残差差异 $<0.1\%$（静态刚度锁定），速度级 $\|e_\xi\|_{\mathrm{rms}}$（$\times10^{-3}$）如下：

\begin{center}
\begin{tabular}{@{}lccc@{}}
\toprule
工况 & C1 & C2 & C3 \\
\midrule
标准 & 0.940 & 0.940 & 0.955 \\
高速 ($\omega=2.5$) & 2.314 & 2.314 & 2.361 \\
噪声 & 3.164 & 3.165 & 3.234 \\
\bottomrule
\end{tabular}
\end{center}

\textbf{近恒等域的结构等价}：C1 与 C2 在全部工况下数值等价（相对差 $\le 0.05\%$），确认二者共享同一信息集（解析前馈 + $\mathrm{Ad}$-搬运 twist 误差）。差异仅在于位姿反馈整形（$A^\top$ vs 螺旋对数），在近恒等域为高阶项——这定量支持了本文的核心主张：C1 相对 C2 的优势不在精度，而在证书（定理 3 的耗散等式、水平集不变性、均方界对 $A^\top$ 整形成立，对螺旋对数不成立）。

\textbf{极限工况的结构性分离}：C3 在噪声下劣化最明显（3.234 vs 3.165，相对 +2.2\%），与其桥接差分 $\Delta\dot q_{\mathrm{cmd}}/\Delta t$ 的噪声放大机制一致；在高速下劣化 +2.0\%，源于其一阶前馈无法解析补偿向心加速度的 $\omega^2$ 项。C1/C2 的 TNDQ 解析前馈避免了数值差分，故在极限工况下保持紧致——这验证了“有严格动力学证书”与“无证书桥接”在工程上的分化。

\subsubsection{水平集余度（验证 V3）}

杯附着后负载突变瞬间，$V^{\mathrm{base}}$ 峰值 $2.47\times10^{-2}$，在 $1.50\,\text{s}$ 内回落至 $3.36\times10^{-4}$；换算至 tuned 档权重后 $V^{\mathrm{tuned}}_{\mathrm{peak}}\le 0.494$，距定理 3(b) 水平集阈值 $c^*=160$ 余度约 \textbf{2.5 个数量级}。这表明定理 3(b) 的局部结论在实际运行中具有充足的工程裕度，并非仅存在于任意小邻域的理论陈述。

\subsection{讨论：理论优势的仿真映射}

上述结果可映射回 §3--5 的理论结构如下：

\textbf{(i) 定理 3(d) 的实用性}：$L_\infty$ 均方界 \eqref{eq:5.7} 预测了偏差型扰动下的稳态误差水平，而 \eqref{eq:5.9} 进一步给出了增益整定的解析规则——仿真中 1.93\% 的反演一致性确认了该规则的有效性。

\textbf{(ii) 定理 3(b) 的工程余度}：水平集 $\Omega_c$ 在实际运行中留有 2.5 个数量级余度，说明局部指数稳定结论在典型工况下是保守但可用的。

\textbf{(iii) 证书分化的可观测性}：C1 与 C2 的数值等价确认了“同信息集”前提，使得“C1 有证书、C2 无证书”成为可辩护的理论分化而非精度吹嘘；C3 在极限工况下的系统性劣化（噪声/高速）则从反面验证了动力学证书的价值——无严格前馈结构时，离散化与差分误差在应力下累积。

\textbf{(iv) 局限}：每组仅单次运行，1.5\%--2\% 的相对差异不具备统计显著性；真机实验与更大负载范围为后续工作。
